\bookchapter{Capitolo 25: Gestione del telefono usato ed eredità digitale}

\chapterintro{Ehi ragazzi, dopo oltre venti capitoli di “addestramento”, dovresti essere cresciuto da “novizio digitale” a un utilizzatore di telefoni abbastanza disinvolto. Ma la vita tecnologica ha i suoi cicli — il nuovo amore diventa vecchio. Quando decidi di “ritirare” il tuo vecchio compagno, sai come salutarlo in modo dignitoso e sicuro?}

\booksection{“Eutanasia” del vecchio telefono: guida alla vendita o donazione sicura}

Principio chiave: la pulizia dei dati non è “eliminazione”, ma “sovrascrittura”. Molti pensano che spostando le foto nel cestino e svuotando l'album sia tutto a posto; poi, dopo la cessione, il nuovo proprietario con un software di recupero dati ritrova vecchie foto, cronologia chat, foto della carta d'identità — queste “figuracce sociali” si ripetono ogni anno.

\begin{enumerate}
  \item \textbf{"Pulizia generale" fai da te}: completa l'ultimo backup integrale, esci da tutti gli account (iCloud, Huawei/Xiaomi/Samsung, Google), disattiva "Trova il mio telefono" — altrimenti il prossimo utente sarà bloccato dal blocco attivazione.
  \item \textbf{"Formattazione" finale}: il "ripristino di fabbrica" dei telefoni moderni di solito include una cancellazione crittografata — distrugge la chiave di cifratura stessa; una volta persa la chiave, tutti i dati sul disco diventano immediatamente illeggibili.
  \item \textbf{Ultimo "trucco di mascheramento" (solo per vecchi dispositivi)}: i telefoni moderni (iOS 8+, Android 6+) hanno già la crittografia completa del disco (FBE) predefinita; il ripristino di fabbrica distrugge la chiave di cifratura, i dati sono teoricamente irreversibili,\textbf{questo passaggio non serve}. Solo per dispositivi molto vecchi non crittografati (anteriore al 2015) può essere necessario, dopo il ripristino di fabbrica, riempire la memoria con file spazzatura e poi eseguire di nuovo il ripristino.
\end{enumerate}

\textbf{Tabella di consultazione rapida dei percorsi di cancellazione dati per brand}: pur essendo simili i principi, ogni brand nasconde il “ripristino di fabbrica” in posizioni leggermente diverse; le differenze chiave sono le seguenti.

\begin{center}
\begin{tabularx}{\linewidth}{|X|X|X|}
\hline
\textbf{Brand} & \textbf{Punto di accesso alla cancellazione} & \textbf{Differenze chiave e avvertenze} \tabularnewline
\hline
Huawei & Impostazioni > Sistema e aggiornamenti > Ripristina > Ripristino dati di fabbrica & Prima uscire dall'account Huawei e disattivare "Trova il mio telefono", altrimenti si innesca il blocco attivazione; in modalità pura può essere richiesta la verifica della password. \tabularnewline
Xiaomi / Redmi & Impostazioni > Il mio dispositivo > Ripristino dati di fabbrica & In fondo ci sono le opzioni "Formatta SD virtuale" e "Cancella tutti i contenuti e le impostazioni": entrambe vanno spuntate; prima di uscire dall'account Xiaomi disattivare "Trova dispositivo". \tabularnewline
OPPO / OnePlus & Impostazioni > Altre impostazioni > Ripristina telefono > Cancella completamente tutti i dati & ColorOS offre una modalità "cancellazione completa" che sovrascrive più volte la memoria; richiede la password di sblocco schermo per la conferma. \tabularnewline
vivo & Impostazioni > Gestione sistema > Backup e ripristino > Cancella tutti i dati & Spuntare "Formatta memoria telefono" per pulire anche l'area SD; uscire dall'account vivo, disattivare "Trova telefono". \tabularnewline
Samsung & Impostazioni > Gestione generale > Ripristina > Ripristino dati di fabbrica & Prima rimuovere l'account Samsung e disattivare "Trova il mio telefono"; alcuni modelli hanno un'opzione di "sovrascrittura automatica", consigliata. \tabularnewline
iPhone & Impostazioni > Generale > Trasferisci o inizializza iPhone > Inizializza tutti i contenuti e le impostazioni & Bisogna prima uscire dall'Apple ID, altrimenti il prossimo utente è bloccato dal blocco attivazione iCloud; iOS 15+ supporta "mantieni dati iCloud all'uscita" per verifiche future. \tabularnewline
\hline
\end{tabularx}
\end{center}

\begin{tcbwarn}{Le tre cose più facilmente trascurate}
\begin{itemize}
  \item \textbf{Non essere usciti dall'account}: dopo la vendita il telefono viene bloccato in remoto, e non riesci più a uscirne; alla fine devi andare con il documento dal servizio ufficiale per sbloccarlo.
  \item \textbf{Non aver rimosso la SD}: il ripristino di fabbrica non formatta le foto sulla SD esterna, ricordati di estrarla o formattarla separatamente.
  \item \textbf{Periferiche USB/OTG non scollegate}: se durante la cancellazione ci sono periferiche collegate, potresti cancellare per sbaglio i loro dati o interrompere la cancellazione.
\end{itemize}

\end{tcbwarn}

\booksection{Eredità digitale: dove collocare la tua anima sociale?}

I dati nel telefono non sono più “virtuali” da tempo. Portano il tuo denaro, le relazioni sociali, le opere che hai creato e persino la tua “identità digitale”.\textbf{Quando una persona muore, come vanno gestiti questi asset?}Non è più una questione astratta, ma un problema reale per sempre famiglie: i parenti non conoscono le password, le piattaforme non permettono l'eredità, le criptovalute sono perse per sempre.

\begin{tcbinfo}{Base giuridica: art. 49 della Legge sulla protezione delle informazioni personali}
L'art. 49 della “Legge della Repubblica Popolare Cinese sulla protezione delle informazioni personali”, in vigore dal 2021, stabilisce chiaramente:\textbf{in caso di morte di una persona fisica, i parenti prossimi, per un loro interesse legittimo e giustificato, possono esercitare sui dati personali del defunto i diritti di consultazione, copia, rettifica e cancellazione previsti in questo capitolo; salvo diversa disposizione lasciata dal defunto in vita.}

Ciò significa:

\begin{itemize}
  \item \textbf{I parenti prossimi hanno un fondamento giuridico}per chiedere alle piattaforme di gestire account e dati personali del defunto.
  \item L'"interesse legittimo e giustificato" di solito include la gestione dell'eredità, la chiusura degli account, l'eliminazione di privacy sensibile.
  \item Se il defunto aveva disposto diversamente con testamento o accordo (ad esempio designando un contatto per l'eredità), si segue la sua volontà.
  \item Le piattaforme non possono rifiutarsi a priori di collaborare in nome della "privacy dell'utente", ma le procedure specifiche variano.
\end{itemize}

\end{tcbinfo}

\booksubsection{WeChat e Alipay: l'account non si eredita, ma il denaro sì}

Queste due app, le più usate dai cinesi, hanno regole diverse:

\begin{itemize}
  \item \textbf{Account WeChat}: secondo la "Licenza e accordo di servizio del software Tencent WeChat", la proprietà dell'account spetta a Tencent, l'utente ha solo il diritto di uso;\textbf{l'account in sé non è ereditabile e non è cedibile}. Ma le\textbf{giacenze di cambio e i fondi di Qianlitong}nell'account sono beni personali: i parenti prossimi possono chiederne il prelievo o il trasferimento a Tenpay con certificato di morte, certificato di parentela e atto notarile di successione.
  \item \textbf{Alipay}: anche l'account non è direttamente ereditabile, ma il saldo Alipay, Yu'ebao, Licaitong e altri\textbf{fondi sono beni ereditabili}. Gli eredi possono rivolgersi all'assistenza Alipay o presentare di persona certificato di morte e atto notarile di successione; i prodotti di investimento di Ant Fortune richiedono la procedura della società di fondi.
  \item \textbf{Cronologia chat}: legalmente è privacy personale, la piattaforma non la fornisce spontaneamente alla famiglia. Se è rilevante per l'eredità, può essere ottenuta tramite cause o richiesta notarile.
\end{itemize}

\begin{tcbwarn}{Cose da non fare mai}
\begin{itemize}
  \item \textbf{Non consegnare direttamente a un familiare le password del defunto "per continuare a usarlo"}— viola l'accordo di piattaforma, e se i fondi vengono congelati o si innesca il controllo rischi, è più difficile fornire prove.
  \item \textbf{Non provare a "prelevare" usando il telefono del defunto per ricevere i codici di verifica}— può essere considerato furto d'uso, con possibili conseguenze penali.
  \item \textbf{Non temporeggiare}: alcuni account, dopo lungo periodo di inattività, vengono cancellati automaticamente; recuperare i fondi diventa più complicato.
\end{itemize}

\end{tcbwarn}

\booksubsection{“Contatto per l'eredità” di Apple: designarlo prima rende tutto più facile}

Da iOS 15.2 Apple ha introdotto la funzione “Contatto per l'eredità”, che permette di designare in anticipo un contatto fidato; dopo il decesso, questi può richiedere l'accesso ai dati iCloud del defunto (foto, note, contatti, memo, ecc.) con il\textbf{certificato di morte + chiave di accesso del contatto per l'eredità}. Passaggi:

\begin{enumerate}
  \item Sull'iPhone vai su\textbf{[Impostazioni] > [il tuo nome] > [Accesso e sicurezza] > [Contatto per l'eredità]}.
  \item Tocca "Aggiungi contatto per l'eredità", conferma l'identità con Face ID.
  \item Scegli il contatto (dalla rubrica), il sistema genererà una\textbf{chiave di accesso}.
  \item Consegna la chiave al contatto tramite SMS, email o stampa cartacea, ricordandogli di custodirla con cura.
  \item Puoi aggiungere più contatti come backup, per evitare un singolo punto di fallimento.
\end{enumerate}

La persona designata non ha bisogno di un dispositivo Apple: con chiave e certificato di morte può richiedere l'accesso sul sito Apple. L'intera procedura di solito si completa in 7 giorni.

\booksubsection{“Gestione account inattivo” di Google: gestire il dopo di noi in automatico}

Google offre l'“Inactive Account Manager”, in grado di eseguire automaticamente operazioni presetstate dopo un lungo periodo di inattività dell'utente — è uno dei meccanismi di eredità digitale più completi attualmente disponibili.

\begin{enumerate}
  \item Visita\textbf{myaccount.google.com/inactive}, accedi al tuo account Google.
  \item Imposta il periodo di timeout: 3 / 6 / 12 / 18 mesi di inattività prima dell'attivazione.
  \item Aggiungi contatti fidati (massimo 10), impostando per ciascuno a quali dati può accedere (Gmail, Foto, Drive, ecc.).
  \item Opzionale: imposta una risposta email automatica che informi chi ti scrive che non puoi più rispondere.
  \item Scegli infine: alla scadenza,\textbf{condividere i dati con i contatti}oppure\textbf{cancellare direttamente l'account}.
  \item Imposta un numero di telefono come verifica di fallback, per evitare attivazioni per errore.
\end{enumerate}

\booksubsection{Account social cinesi: commemorazione e cancellazione}

Le principali piattaforme social cinesi trattano gli account dei defunti in due modi:\textbf{“commemorazione” (mantenere l'account per il ricordo dei familiari)}e\textbf{“cancellazione” (eliminazione completa)}. I familiari possono fare richiesta all'assistenza della piattaforma con certificato di morte, certificato di parentela e documento d'identità del richiedente.

\begin{itemize}
  \item \textbf{Weibo}: supporta la "commemorazione" dell'account, su richiesta di un familiare con certificato di morte; l'account viene bloccato, non è più accessibile ma i contenuti restano, e sulla home compare "Questo account è stato impostato come account commemorativo". Si può anche chiedere la cancellazione, fornendo un atto notarile di successione.
  \item \textbf{Douyin (TikTok Cina)}: tramite "Io - Impostazioni - Feedback e aiuto - Account e sicurezza - Disattivazione account" o via assistenza, si presentano certificato di morte, certificato di parentela e documento del richiedente; dopo l'approvazione si può disattivare o trasformare in account commemorativo.
  \item \textbf{Bilibili}: tramite assistenza si presentano certificato del defunto e certificato di parentela; si può chiedere la disattivazione o la conservazione come pagina commemorativa; il trattamento dei contenuti caricati dall'UP deve tener conto delle dichiarazioni di copyright.
  \item \textbf{Xiaohongshu}: tramite assistenza ufficiale, con certificato di morte e certificato di parentela, si può richiedere la disattivazione dell'account e la cancellazione dei contenuti.
\end{itemize}

\begin{tcbinfo}{Promemoria chiave}
\begin{itemize}
  \item Le regole delle piattaforme si aggiornano: prima di procedere consulta le ultime indicazioni dell'assistenza ufficiale.
  \item Si consiglia di scansionare e fare backup di tutti i documenti prima di consegnarli, conservando con cura gli originali.
  \item Per procedure una tantan su più piattaforme, si può richiedere un unico\textbf{atto notarile di successione}, evitando di ripetere più volte le pratiche.
\end{itemize}

\end{tcbinfo}

\booksubsection{La complessa successione degli asset digitali}

Oltre agli account social, sempre più persone detengono\textbf{asset digitali di vero valore}, la cui successione è più complessa:

\begin{itemize}
  \item \textbf{Criptovalute (Bitcoin, Ethereum, ecc.)}: sul principio "la chiave privata è la proprietà",\textbf{senza chiave privata l'asset è perso}. Gli eredi non possono "ricorrere" on-chain con documenti legali. Consiglio: conservare la seed phrase/chiave privata in frammenti in più posizioni fisiche sicure (come cassetta di sicurezza in banca + cassaforte di casa) e indicare nel testamento le modalità di accesso.
  \item \textbf{Account sugli exchange}: come i conti bancari, si può chiedere agli exchange (come Binance, OKX) la gestione con atto notarile di successione, ma le procedure su exchange cross-border sono lunghe.
  \item \textbf{Domini}: sono beni ereditabili; gli eredi possono trasferirli presso il registrar (come Alibaba Cloud, GoDaddy) presentando l'atto notarile di successione.
  \item \textbf{Diritti di proprietà intellettuale (articoli originali, video, brevetti, diritti d'autore del software)}: i diritti patrimoniali del diritto d'autore sono ereditabili (i diritti morali no); il diritto ai proventi spetta agli eredi fino a 50 anni dopo la morte dell'autore.
  \item \textbf{Cloud storage e netdisk}: i dati personali sono ereditabili; ma alcuni accordi di piattaforma prevedono che l'account non sia cedibile: occorre prima scaricare e poi disattivare.
\end{itemize}

\textbf{Il vero costo della “scomparsa” delle criptovalute}\par

Secondo un rapporto Chainalysis, nel mondo sono già andati persi per sempre\textbf{oltre 3 milioni di Bitcoin}a causa di chiavi private smarrite, pari a circa un sesto del totale esistente. Se detieni criptoattività, assicurati che almeno una persona di assoluta fiducia conosca il\textbf{percorso di recupero}— altrimenti queste ricchezze resteranno per sempre addormentate sulla blockchain.

\booksubsection{Un modello di “testamento digitale”}

Il testamento patrimoniale tradizionale non copre gli asset digitali. Si consiglia di aggiornare ogni anno un allegato “testamento digitale”, da conservare insieme al testamento ufficiale. Modello di riferimento:

\textbf{“Allegato per la gestione degli asset digitali personali” (modello)}\par

\textbf{Parte 1: elenco account}(registra solo nome utente e piattaforma,\textbf{non le password})

\begin{itemize}
  \item Social: WeChat (account xxx), Weibo, Douyin, Bilibili, Xiaohongshu, ecc.
  \item Pagamenti: Alipay, WeChat Pay, app delle varie banche
  \item Cloud: iCloud, Huawei Cloud, Mi Cloud, account Google, account Microsoft
  \item Email: QQ Mail, NetEase Mail, Gmail, ecc.
  \item Altro: netdisk, servizi in abbonamento, account di gioco, ecc.
\end{itemize}

\textbf{Parte 2: soluzione di custodia delle password}

\begin{itemize}
  \item Custodia della password principale: usa un password manager (come Bitwarden, 1Password); la password principale è scritta su carta sigillata nella cassaforte di casa.
  \item Chiavi di recupero/seed phrase: custodite a frammenti, almeno in due posizioni fisiche (come cassetta in banca + cassaforte di casa).
  \item Persona autorizzata alla consultazione: nome, documento d'identità, rapporto con l'interessato, contatto.
\end{itemize}

\textbf{Parte 3: volontà di gestione}

\begin{itemize}
  \item Per ciascun account, preferenza tra "commemorazione" e "cancellazione".
  \item Se per iCloud / Google dati si designa un "contatto per l'eredità" o si imposta l'"Inactive Account Manager".
  \item Chi eredita le criptoattività, con spiegazione del percorso per ottenere la chiave privata.
  \item A chi spettano i proventi delle opere coperte da diritto d'autore.
\end{itemize}

\textbf{Parte 4: contatti e piattaforme}

\begin{itemize}
  \item Esecutore fiduciario: nome, contatto, contatto di backup.
  \item Elenco degli accessi all'assistenza delle piattaforme e dei documenti necessari.
  \item Data di aggiornamento dell'allegato: \_\_\_\_/\_\_/\_\_; prossima revisione: \_\_\_\_/\_\_/\_\_.
\end{itemize}

\begin{tcbwarn}{Nota legale sulle "password"}
Scrivere le password direttamente nel testamento comporta rischi di fughe di informazioni; una volta pubblicato, tutti gli eredi potrebbero vederle.\textbf{È più prudente: nel testamento inserire solo “elenco account + volontà + luogo di custodia delle password”, e trasmettere la vera password — tramite la master password del password manager — solo all'esecutore fiduciario}. Così si rispettano le esigenze legali e si riduce il rischio di fughe.

\end{tcbwarn}

\booksection{L'ultimo viaggio ecosostenibile}

Un telefono contiene decine di metalli tra cui oro, argento, rame e terre rare; disfarsene a caso inquina il suolo e spreca risorse. I canali di riciclo corretti sono tre, con punti di forza diversi:

\begin{center}
\begin{tabularx}{\linewidth}{|X|X|X|X|}
\hline
\textbf{Canale} & \textbf{Prezzo di recupero} & \textbf{Sicurezza dati} & \textbf{Scenario d'uso} \tabularnewline
\hline
Permuta ufficiale (Huawei/Apple/Xiaomi, ecc.) & Medio, spesso con buoni integrati & Massima, con procedure ufficiali & Più conveniente quando si passa a un nuovo telefono dello stesso brand \tabularnewline
Piattaforma di riciclo di terze parti (Aihuishou, Zhuanzhuan, Huishoubao, ecc.) & Leggermente super, confrontabile & Alta, con report di cancellazione dati & Per chi vuole monetizzare direttamente, senza vincoli di brand \tabularnewline
Donazione di beneficenza (scuole di campagna, enti benefici) & Nessun ritorno economico & Richiede cancellazione autonoma completa & Dispositivo ancora funzionante, voglia di fare del bene \tabularnewline
Riutilizzo (ritenuto e modificato) & Risparmio su un nuovo dispositivo & Continua a essere usato dal proprietario originale, nessun rischio & Dispositivo ancora performante, esigenze di riuso in casa \tabularnewline
\hline
\end{tabularx}
\end{center}

\begin{itemize}
  \item \textbf{Riciclo/permuta ufficiale}: la soluzione più consigliata, procedure standard e sicurezza dei dati garantita. I siti di Apple, Huawei, Xiaomi hanno sezioni "permuta" dove, dopo valutazione, si ottiene uno sconto sul nuovo telefono.
  \item \textbf{Piattaforme di riciclo regolari}: scegli grandi piattaforme con buona reputazione; dopo l'ordine, ritiro a domicilio e pagamento dopo verifica. Scegli sempre piattaforme che\textbf{forniscono un "report di cancellazione dati"}e conservane la ricevuta.
  \item \textbf{Donazione di beneficenza}: puoi contattare scuole di campagna locali o organizzazioni no profit; i vecchi telefoni possono essere riadattati come tablet didattici o telefoni per anziani.\textbf{Prima della donazione esegui sempre la procedura completa di cancellazione e rimuovi SIM/SD}.
  \item \textbf{Riutilizzo}: trasformalo in navigatore auto, centro di controllo domotico, telecamera di sorveglianza, baby monitor, cornice digitale. Un vecchio telefono può ancora rendere per anni.
\end{itemize}

\begin{tcbinfo}{La verità sui rifiuti elettronici}
Il circuito di un telefono contiene metalli preziosi come\textbf{oro, argento, palladio, rame}; ogni tonnellata di circuiti usati contiene decine di volte più oro di un minerale aurifero. Ma se smaltito in modo scorretto, il piombo, il mercurio e il cadmio penetrano nel suolo e nelle falde, inquinando per decenni.\textbf{Un riciclo corretto protegge l'ambiente e rimette in circolo metalli scarsi.}

\end{tcbinfo}

\begin{tcbpractice}{Pratica: un "funerale standard" per il tuo vecchio telefono}
\begin{enumerate}
  \item Ultimo backup completo: carica le foto sul cloud, copia i file importanti sul computer, esporta la cronologia chat.
  \item Esci da tutti gli account: iCloud / Huawei / Xiaomi / Google, disattiva "Trova il mio telefono".
  \item Estrai SIM e SD, scollega tutte le periferiche OTG.
  \item Esegui il ripristino di fabbrica con cancellazione crittografata, se necessario spuntando manualmente "formatta memoria".
  \item (Solo vecchi dispositivi) Riempi la memoria con file spazzatura e fai di nuovo il ripristino di fabbrica.
  \item Scegli la sua prossima tappa: permuta ufficiale, piattaforma di riciclo regolare, donazione di beneficenza o donazione a un familiare.
  \item Aggiorna il tuo "elenco eredità digitale", cancellando i dispositivi dismessi.
\end{enumerate}

\end{tcbpractice}

\begin{tcbpractice}{Lista delle mansioni del mese}
\begin{enumerate}
  \item Per il tuo telefono principale, designa un "contatto per l'eredità" Apple o imposta l'"Inactive Account Manager" di Google.
  \item Usa il modello di questa sezione per redigere un "Allegato per la gestione degli asset digitali personali" e accordati con la famiglia sul luogo di conservazione.
  \item Verifica se in casa ci sono 2 o più vecchi telefoni inutilizzati e prenota un ritiro ufficiale o su piattaforma regolare.
\end{enumerate}

\end{tcbpractice}
